\chapter{What the failures had in common}

\section{The dangerous defects were the silent ones}

Of the eleven entries in this report, three crashed, three failed to compile,
and five produced plausible output while being wrong. The five silent ones took
almost all the time.

A build targeting Turing instead of Blackwell compiles, links and runs. A
benchmark whose columns have shifted parses cleanly and plots. An optimisation
that relieves a resource nothing was waiting on delivers a small positive
number that is easy to accept. A neighbour search that counts some particles
twice produces a fluid that still behaves like a fluid. A diagnostic that
perturbs the state it measures produces a smooth curve.

None of these announces itself, and none is found by testing that the program
runs. Each was caught by a different mechanism, and the pattern is that each
mechanism compares the thing against something independent: the object file
against the intended architecture, the figure against what a tile width could
plausibly be, the optimisation against a bound computed from measured
bandwidth, the neighbour search against a brute force sum, the diagnostic run
against an undiagnosed one.

\section{A tolerance without a scale is not a tolerance}

Two defects came from this directly. The kernel continuity test compared a
quantity of order $10^6$ against an absolute tolerance of $10^{-3}$ and failed a
correct implementation. Had the comparison gone the other way, against a
quantity of order $10^{-6}$, it would have passed anything at all.

Every gate in the project is now expressed against a scale that means
something. The momentum gate is bounded by $\epsilon_{\mathrm{f32}} \sqrt{N}
n_{\mathrm{steps}}$, which is what single precision arithmetic can explain, and
it passes with six orders of magnitude to spare, so the number reported is
informative rather than merely a pass. The force gate is relative to the
magnitude of the acceleration being checked. The density gate is relative to
rest density.

\section{Size the optimisation before writing it}

The fusion result is the clearest lesson in the project. The saving was
computable in advance from two numbers already measured at Phase 0: the bytes
removed and the achievable bandwidth. That calculation gives
\FusionCeilingPercent{} percent as an absolute ceiling at \LadderMaxN{}
particles and takes two minutes.

The same calculation would have reframed the tiling work correctly as well.
Tiling is worth a factor of \TilingSpeedupAtMaxN{}, which no bandwidth argument
can explain, because bandwidth was never the constraint. Knowing that in
advance would have pointed at the load store pipeline from the beginning
instead of at the memory bus.

\section{Verification has to be independent of the thing verified}

The strongest checks in this project are the ones where the reference was
produced by a genuinely different route. The double precision host force
implementation shares no code with the kernels and is written as literally as
the force law can be written. The brute force density sum shares no code with
the neighbour search. The chi squared test is checked by feeding it a
deliberately wrong model that it must reject, so a test that could never fail
would itself fail.

Where sharing is unavoidable it is deliberate and in the safe direction: the
smoothing kernels and the force law are compiled from one header into both the
device kernels and the host tests, so the two cannot drift apart. That is
sharing an implementation to guarantee agreement, not sharing an assumption to
guarantee a pass.

\section{The environment is part of the experiment}

Four of the eleven entries were toolchain problems, and one of them would have
invalidated every performance number in the main report while leaving no trace
except numbers that were slightly disappointing. Phase 0 of this build existed
to probe the environment and record what it found, and that turned out to be
time well spent rather than ceremony. The measured bandwidth figure recorded
there is the denominator of the fusion ceiling calculation that explains the
central performance result of the project.
